\newcommand{\rb}[1]{\left(#1\right)}
\newcommand{\bb}[1]{\left[#1\right]}
\newcommand{\cb}[1]{\left\{#1\right\}}
\newcommand{\floor}[1]{\left\lfloor#1\right\rfloor}
\newcommand{\ceil}[1]{\left\lceil#1\right\rceil}
\newcommand{\round}[1]{\left\lfloor#1\right\rceil}
\newcommand{\floorfrac}[2]{\floor{\frac{#1}{#2}}}
\newcommand{\cf}[2]{#1\!\rb{#2}}
\newcommand{\ctf}[2]{\textsc{#1}\!\rb{#2}}
\newcommand{\fmax}[1]{\ctf{Max}{#1}}
\newcommand{\fmin}[1]{\ctf{Min}{#1}}
\newcommand{\ftan}[1]{\ctf{tan}{#1}}
\newcommand{\fsin}[1]{\ctf{sin}{#1}}
\newcommand{\fcos}[1]{\ctf{cos}{#1}}
\newcommand{\fasin}[1]{\ctf{asin}{#1}}
\newcommand{\facos}[1]{\ctf{acos}{#1}}
\newcommand{\abs}[1]{\left|#1\right|}
\newcommand{\Set}[2]{\{\,#1\mid#2\,\}}
\newcommand{\hide}[1]{\iffalse #1 \fi}
\newcommand{\norm}[1]{\left\|\,#1\,\right\|}
\newcommand{\enorm}[1]{\norm{#1}_2}
\newcommand{\nnat}{\mathbb{N}}
\newcommand{\nint}{\mathbb{Z}}
\newcommand{\nrat}{\mathbb{Q}}
\newcommand{\nreal}{\mathbb{R}}
\newcommand{\ncomp}{\mathbb{C}}

\newcommand{\nrealop}{\mathbb{R}_{\geq 0}}
\newcommand{\nrealp}{\mathbb{R}_{>0}}

\DeclareOldFontCommand{\rm}{\normalfont\rmfamily}{\mathrm}
\DeclareOldFontCommand{\sf}{\normalfont\sffamily}{\mathsf}
\DeclareOldFontCommand{\tt}{\normalfont\ttfamily}{\mathtt}
\DeclareOldFontCommand{\bf}{\normalfont\bfseries}{\mathbf}
\DeclareOldFontCommand{\it}{\normalfont\itshape}{\mathit}
\DeclareOldFontCommand{\sl}{\normalfont\slshape}{\@nomath\sl}
\DeclareOldFontCommand{\sc}{\normalfont\scshape}{\@nomath\sc}

\newcommand{\nvec}[1]{\hat{#1}}
\newcommand{\mat}[1]{\mathbf{#1}}
\newcommand{\quat}[1]{\underline{#1}}
\newcommand{\pnt}[1]{\mathbf{#1}}

% Where the supplemental material is: the separate supplemental.tex, or the appendix of the
% arXiv version, whose main file (make arxiv) defines this name first.
\providecommand{\SupplementalName}{supplemental material}
\ifdefined\ArxivVersion% Settings of the arXiv version (make arxiv); loaded from common/commands.tex only if
% \ArxivVersion is defined.
%
% The arXiv version appends the supplemental material to the paper as appendix. If the paper
% already has an appendix, the second \appendix continues its lettering instead of starting
% again at A.
\let\ArxivOriginalAppendix\appendix
\newif\ifArxivInAppendix
\def\appendix{%
  \ifArxivInAppendix
    \edef\ArxivSectionCount{\the\value{section}}%
    \ArxivOriginalAppendix
    \setcounter{section}{\ArxivSectionCount}%
  \else
    \ArxivOriginalAppendix
  \fi
  \global\ArxivInAppendixtrue}
% Overrides of the document class, e.g., to remove conference markings, if the template has any.
\IfFileExists{template/arxiv.tex}{\input{template/arxiv.tex}}{}
\fi

% \todo{who or what}{text}: a highlighted note in the pdf.
\definecolor{todocolor}{RGB}{220,0,120}
\newcommand{\todo}[2]{%
  {\color{red}\textbf{[TODO: #1] }}%
  {\color{todocolor}{#2}}%
}
